\section{Тема 10. Алгоритмы обработки символьной информации}

\subsection{Постановка задачи (вариант~4)}

В заданном тексте найти самое длинное слово и самую длинную фразу.
Фразой считается последовательность символов между знаками
завершения предложения (\texttt{.}, \texttt{!}, \texttt{?}).

\subsection{Математическое обоснование}

Текст обрабатывается за один линейный проход. Слова выделяются
как максимальные подпоследовательности буквенных символов;
по ходу прохода ведётся текущая длина слова~$w$ и максимальная
длина $w_{\max}$. Аналогично для фраз~--- ведётся текущая
длина~$p$ и максимальная~$p_{\max}$, со сбросом при встрече
разделителя предложения. Сложность~--- $O(L)$, где $L$~--- длина
текста.

\subsection{Блок-схема алгоритма}

\begin{center}
\begin{tikzpicture}[
  node distance=10mm, start chain=going below,
  nodes={draw=black, top color=white, bottom color=lightgray!50,
         minimum width=5.6cm, align=center, font=\small},
  arrow/.style={->, >=Stealth, thick},
]
  \node[draw, rounded corners, on chain] {Start};
  \node[shape=trapezium, trapezium left angle=70,
        trapezium right angle=110, on chain, draw, join=by arrow]
       {$T$};
  \node[shape=chamfered rectangle, chamfered rectangle xsep=10pt,
        on chain, draw, join=by arrow]
       (loop) {$i := 0,\ |T|$};
  \begin{scope}[local bounding box=scl]
    \node[shape=diamond, aspect=2, on chain, draw, join=by arrow,
          minimum width=5.0cm]
         (c) {$T[i]$ -- буква?};
    \node[shape=rectangle, draw, right=2cm of c]
         (wa) {$w \mathrel{+}= 1$};
    \draw[arrow] (c.east) -- node[above,font=\scriptsize]{$+$} (wa.west);
    \node[shape=rectangle, on chain, draw, below=14mm of c]
         (wu) {$w_{\max}$ обновить,\ $w:=0$};
    \draw[arrow] (c.south) -- (wu.north);
    \draw[arrow] (wa.south) |- (wu.east);
    \node[shape=rectangle, on chain, draw, join=by arrow]
         {$p_{\max}$ обновить при \texttt{.!?}};
    \coordinate[on chain, join=by arrow] (le);
  \end{scope}
  \draw[arrow] (scl.south) -| ($(scl.west)-(0.9,0)$) |- (loop);
  \coordinate[on chain] (sp);
  \draw[arrow] (loop) -| ($(scl.east)+(0.9,0)$) |- (sp);
  \node[shape=trapezium, trapezium left angle=70,
        trapezium right angle=110, on chain, draw, join=by arrow]
       {$w_{\max}$,\ $p_{\max}$};
  \node[draw, rounded corners, on chain, join=by arrow] {End};
\end{tikzpicture}
\end{center}

\subsection{Текст программы}

\lstinputlisting[style=Cstyle]{../src/t10.c}

\subsection{Тестовые данные}

\textbf{Результат выполнения программы:}
\begin{lstlisting}[language={},basicstyle=\ttfamily\small,frame=single]
Введите текст (одна строка):
Программирование является увлекательным занятием. Алгоритмы важны!
Самое длинное слово: "Программирование" (17 букв)
Самая длинная фраза: " Программирование является увлекательным занятием"
\end{lstlisting}
